\ExplSyntaxOff

\RpgIfSections{
	%%%%%%%%%%%%%%%%%%%%%%%%%%%%%%%%%%%%%%%%%
	% Section formatting
	%%%%%%%%%%%%%%%%%%%%%%%%%%%%%%%%%%%%%%%%%
	% Part
	%% set explicitly, so that switching back from a theme which changes it (such as scifi) restores Roman numerals
	\renewcommand{\thepart}{\Roman{part}}
}
\DeclareDocumentCommand\RpgThemeSetPHBParts{}
{
	\DeclareDocumentCommand\RpgThemePartPageTikz{}
	{
		\def\radius{0.18\paperwidth}
		\def\overset{7}
		\begin{tikzpicture}
			\draw[line width=3pt,TitleGold,fill=BrTan] (0,0) ++ ({-\radius*cos(\overset)},{8pt}) arc ({180-\overset}:{360+\overset}:\radius)--cycle;
		\end{tikzpicture}
	}
}
\DeclareDocumentCommand\RpgThemeSetDMGParts{}
{
	\DeclareDocumentCommand\RpgThemePartPageTikz{}
	{
		\def\w{0.18\paperwidth}
		\def\h{0.12\paperwidth}
		\def\p{0.08\paperwidth}
		\def\f{0.3}
		\begin{tikzpicture}
			\draw[line width=3pt,ContourGray!30!black,fill=BrTan] (-\w,0)--(\w,0)--(\w,-\h)--(0,{-(\h + \p)})--(-\w,-\h)--cycle;

			\node (A) at ({-\w+\f*\w},{-\h - \f*\p}){};
			\node (B) at ({\w-\f*\w},{-\h - \f*\p}){};
			\draw[line width=3pt,ContourGray!30!black] (A) to[out=-20,in=-160] (B);
			\draw[line width=2.1pt,ContourGray] (A) to[out=-20,in=-160] (B);

			\draw[line width=2.1pt,ContourGray] (-\w,0)--(\w,0)--(\w,-\h)--(0,{-(\h + \p)})--(-\w,-\h)--cycle;
		\end{tikzpicture}
	}
}

\RpgThemeSetPHBParts{}
\RpgThemePartPage
{
	\RpgPut[\paperwidth](t){0.5\paperwidth}{\paperheight+0.15cm}{
		\RpgThemePartPageTikz
	}

	\RpgPut[0.4\paperwidth](t){0.5 \paperwidth}{\paperheight - 0.3cm}{
		\RpgFontPart{} \MakeUppercase{Part \Roman{part}}
	}
	\RpgPut[0.4\paperwidth](t){0.5 \paperwidth}{\paperheight - 2.2cm}{
		\RpgFontPartSubtitle{} #1
	}
}


\usetikzlibrary {shadows}
\ExplSyntaxOn

\RpgThemeTitleCover
{
	\def\w{0.3\paperwidth}
	\def\h{0.005\paperwidth}
	\RpgPut[0.9\paperwidth](t){0.5 \paperwidth}{\paperheight-3cm}{ \RpgContour[inner=white,outer=black,style=\RpgFontTitle]{\RpgGetTitle*}
		% %%fancy red diamond
		\begin{center}
			\begin{tikzpicture}
				\fill[red,circular~ glow={fill=black!80!red}] (0,0)--++(\w,-\h)--++(-\w,-\h)--++(-\w,\h)--cycle;
			\end{tikzpicture}
		\end{center}
	}

	\RpgPut[\paperwidth](b){0.5\paperwidth}{1cm}{
		\begin{center}
			\RpgContour[inner=white,outer=black,style={\RpgFontSubTitle}]{\RpgGetSubtitle*}
		\end{center}

		\vspace{1cm}

		{\RpgContour[inner=white,outer=black,style={\RpgFontSubTitle\huge}]{\RpgGetAuthor*}}
	}
}

\ExplSyntaxOff
